\begin{tablalibro}{3}{L{0.28}L{0.34}L{0.38}}
\caption{Estrategias del pensamiento computacional}\label{tab:estrategias-del-pensamiento-computacional}\\
\toprule
\rowcolor{Papel}
\CabeceraTabla{Estrategia} & \CabeceraTabla{Pregunta orientadora} & \CabeceraTabla{Ejemplo} \\
\midrule
\endfirsthead
\SigueTabla{3}
\toprule
\rowcolor{Papel}
\CabeceraTabla{Estrategia} & \CabeceraTabla{Pregunta orientadora} & \CabeceraTabla{Ejemplo} \\
\midrule
\endhead
\midrule
\endfoot
\bottomrule
\multicolumn{3}{@{}p{\linewidth}@{}}{\NotaLibro{Elaboración propia.}}\\
\endlastfoot

Descomposición & ¿En qué partes puedo dividir el problema? & Separar registro, cálculo y reporte de notas. \\
Reconocimiento de patrones & ¿Qué se repite o se parece? & Identificar cálculos iguales en distintos cursos. \\
Abstracción & ¿Qué información necesito conservar? & Considerar las calificaciones y sus pesos. \\
Generalización & ¿Cómo reutilizar la solución? & Aplicar el cálculo a distintos estudiantes. \\
Diseño algorítmico & ¿Qué pasos permiten obtener el resultado? & Validar datos, calcular y comunicar el promedio. \\
\end{tablalibro}
